\documentclass[11pt,a4paper]{article}

\usepackage[a4paper,margin=25mm,headheight=15pt]{geometry}
\usepackage[T1]{fontenc}
\usepackage{lmodern}
\usepackage{microtype}
\usepackage{graphicx}
\usepackage{xcolor}
\usepackage{booktabs}
\usepackage{tabularx}
\usepackage{longtable}
\usepackage{array}
\usepackage{siunitx}
\usepackage{caption}
\usepackage{subcaption}
\usepackage{float}
\usepackage{enumitem}
\usepackage{fancyhdr}
\usepackage{hyperref}
\usepackage[nameinlink,noabbrev]{cleveref}

\definecolor{uomblue}{HTML}{17365D}
\definecolor{softblue}{HTML}{EAF0F7}
\hypersetup{colorlinks=true,linkcolor=uomblue,citecolor=uomblue,urlcolor=uomblue,
  pdftitle={TF Flowers Classification with Convolutional Neural Networks},
  pdfauthor={EN3150 Assignment 03 Group}}
\sisetup{round-mode=places,round-precision=3,group-separator={,}}
\setlength{\parskip}{0.55em}
\setlength{\parindent}{0pt}
\setlength{\emergencystretch}{2em}
\setlist{nosep,leftmargin=1.5em}
\captionsetup{font=small,labelfont=bf}

\pagestyle{fancy}
\fancyhf{}
\fancyhead[L]{\small EN3150 \textbar{} Assignment 03}
\fancyhead[R]{\small TF Flowers Classification}
\fancyfoot[C]{\thepage}
\renewcommand{\headrulewidth}{0.4pt}

\newcommand{\class}[1]{\texttt{#1}}
\newcommand{\reportfigure}[3]{%
  \begin{figure}[H]\centering
  \includegraphics[width=#1\linewidth]{#2}
  \caption{#3}
  \end{figure}}

\begin{document}

% ---------------------------------------------------------------- cover
\hypersetup{pageanchor=false}
\begin{titlepage}
  \centering
  \vspace*{8mm}
  \includegraphics[height=39mm,keepaspectratio]{assets/uom-emblem.jpg}\par
  \vspace{7mm}
  {\Large\bfseries\color{uomblue} UNIVERSITY OF MORATUWA\par}
  \vspace{2mm}
  {\large Department of Electronic and Telecommunication Engineering\par}
  \vspace{13mm}
  {\large EN3150 -- Pattern Recognition\par}
  \vspace{4mm}
  {\large Assignment 03\par}
  \vspace{13mm}
  {\Huge\bfseries\color{uomblue} Flower Image Classification\par}
  \vspace{4mm}
  {\LARGE Convolutional Neural Networks and Transfer Learning\par}
  \vfill
  \begin{tabular}{>{\bfseries}r l}
    Submission type: & Group report \\[4mm]
    Group members: & \begin{tabular}[t]{@{}l@{}}
      230145E -- Dilhan W.A.\\
      230375H -- Lenmini B.L.W.\\
      230070T -- Balasooriya B.R.B.D.\\
      230211E -- Gunsekara K.S.
    \end{tabular} \\[4mm]
    Module: & EN3150 -- Pattern Recognition \\[2mm]
    Dataset: & TF Flowers (3,670 RGB images) \\
  \end{tabular}
  \vfill
  {\large October 2026\par}
\end{titlepage}

\pagenumbering{roman}
\hypersetup{pageanchor=true}
\tableofcontents
\listoffigures
\clearpage
\pagenumbering{arabic}

% ---------------------------------------------------------------- abstract
\begin{abstract}
This project studies five-class flower recognition on the TF Flowers dataset.
The work prepares a reproducible stratified data split, compares optimizers
using a compact depthwise separable convolutional neural network, evaluates a
standard CNN against a parameter-efficient CNN, and fine-tunes two
ImageNet-pretrained mobile architectures~\cite{mobilenet,efficientnet}. The saved experiment records show
that the standard custom CNN achieved 79.82\% test accuracy, while the compact
custom model achieved 76.18\%. MobileNetV2 reached the highest recorded test
accuracy (85.09\%), at substantially higher model size and computational cost
than the compact model. These results illustrate the accuracy--efficiency
trade-off between training a small model from scratch and adapting pretrained
visual features.
\end{abstract}

\textbf{Keywords:} image classification, convolutional neural network,
TF Flowers, depthwise separable convolution, transfer learning

\section{Introduction}
Image classification assigns an input image to one label from a fixed set of
classes. Flower recognition is a useful small-scale vision task because
different species can have similar colour and texture, while examples within a
class vary in pose, scale, background, and lighting. A convolutional neural
network (CNN) learns spatial filters that can represent increasingly complex
visual patterns.

This report documents the implementation and recorded results for five flower
classes: dandelion, daisy, tulips, sunflowers, and roses. It covers dataset
preparation, optimizer selection, custom architecture design, held-out test
evaluation, and fine-tuning of MobileNetV2 and EfficientNet-B0. All numerical
results below are taken from the repository's committed result files; training
was not repeated while preparing this report.

\section{Dataset and experimental design}
\subsection{Dataset}
The project uses the TF Flowers dataset distributed through TensorFlow
Datasets~\cite{tfdsflowers}, containing 3,670 labelled RGB images in five classes. Images are
directly resized to $64\times64$ pixels by bilinear interpolation with
antialiasing. Direct resizing is simple and gives a fixed input shape, but it
can change the original aspect ratio.

\subsection{Reproducible data split}
The preparation script uses seed 42 and performs a separate random shuffle
within each class before assigning approximately 70\% to training, 15\% to
validation, and the remainder to testing. The script checks that every source
image occurs exactly once in the combined split. Integer rounding yields the
counts in \cref{tab:split}; the saved split manifest records the assignment of
each source example.

\begin{table}[H]\centering
\caption{Class counts in the prepared data splits.}\label{tab:split}
\begin{tabular}{lrrrr}
\toprule
Class & Training & Validation & Test & Total \\
\midrule
Dandelion & 629 & 135 & 134 & 898 \\
Daisy & 443 & 95 & 95 & 633 \\
Tulips & 559 & 120 & 120 & 799 \\
Sunflowers & 489 & 105 & 105 & 699 \\
Roses & 449 & 96 & 96 & 641 \\
\midrule
Total & 2,569 & 551 & 550 & 3,670 \\
\bottomrule
\end{tabular}
\end{table}

Only training data is shuffled for model fitting. Custom CNN inputs are
converted to float values in $[0,1]$. Pretrained-model inputs remain in
$[0,255]$ until the architecture-specific preprocessing layer is applied.
The validation set guides model selection. The test set is reserved for final
evaluation in the Q4 and Q5 scripts.

\subsection{Training-only augmentation}
For custom-model and pretrained-model training, the repository applies random
horizontal flips, rotations with a factor of 0.1, and zoom with a factor of
0.1 to training batches only. In Keras, the 0.1 rotation factor corresponds to
up to approximately $\pm36^\circ$; this follows the implementation in
\class{train\_eval.py}, whose descriptive docstring summarizes the rotation
more loosely. Validation and test examples are not augmented.

\section{Methods}
\subsection{Custom CNN architectures}
Both custom models accept $64\times64\times3$ inputs and use four feature
blocks with output channel counts 32, 64, 128, and 256. Each block ends with
$2\times2$ max pooling, reducing the spatial resolution from 64 to 4 pixels.
Global average pooling feeds dropout with rate 0.3, a 128-unit ReLU dense
layer, and a five-unit softmax classifier.

\textbf{Model A (standard CNN)} uses a $3\times3$ Conv2D layer, batch
normalization, and ReLU in each block. Convolution biases are disabled because
the following batch-normalization layer provides a learned offset.

\textbf{Model B (depthwise separable CNN)} replaces each standard convolution
with a $3\times3$ depthwise convolution followed by a $1\times1$ pointwise
convolution. Each operation is followed by batch normalization and ReLU6. This
factorization reduces multiply--accumulate operations (MACs) and parameter
count while retaining the same overall block widths and classifier head.

\subsection{Optimizer comparison}
The optimizer experiment trains Model B for 20 epochs with batch size 32 and
sparse categorical cross-entropy. Runs use the same seed, initial weights,
data split, and shuffled batch order. The compared optimizers are Adam at
$10^{-3}$, SGD at $10^{-2}$, SGD with momentum 0.9 at $10^{-2}$, and a momentum
sweep at 0.5 and 0.99. The validation set is used to compare runs; the saved
script keeps the test split unused for this step.

\subsection{Custom-model training and selection}
The Q4 experiment trains each custom model for 30 epochs with Adam, learning
rate $10^{-3}$, batch size 32, and sparse categorical cross-entropy. Training
uses the same split and data augmentation. A ReduceLROnPlateau schedule is
enabled after the initial epochs, with a minimum learning rate of
$10^{-5}$; checkpoints retain the weights with the best validation accuracy.
The test metrics are computed after selection. The recorded comparison
reports accuracy, macro precision, macro recall, parameter counts, MACs, model
size, and median epoch time.

\subsection{Transfer learning}
MobileNetV2 and EfficientNet-B0 are initialized with ImageNet weights. The
$64\times64$ inputs are bilinearly upsampled to $128\times128$ inside each
model. Training is two-stage: ten epochs train a new five-class head while the
backbone is frozen, then ten epochs fine-tune the top 40 backbone layers at
learning rate $10^{-5}$. The head-only learning rate is $10^{-3}$. Batch
normalization layers in the pretrained backbone remain in inference mode.
The best validation-accuracy checkpoint across both stages is used for final
test evaluation.

\section{Results}
\subsection{Optimizer selection}
\Cref{tab:optimizers} reports the best validation accuracy for each saved run.
Adam reached the highest value, 73.68\% at epoch 20, and was selected for the
custom-model comparison. SGD with momentum 0.9 was the strongest SGD variant
at 72.41\%. The additional momentum sweep did not outperform Adam.

\begin{table}[H]\centering
\caption{Optimizer comparison on Model B (validation set).}\label{tab:optimizers}
\small
\begin{tabular}{lrrr}
\toprule
Optimizer & Learning rate & Best validation accuracy & Best epoch \\
\midrule
Adam & $10^{-3}$ & 0.7368 & 20 \\
SGD & $10^{-2}$ & 0.6987 & 16 \\
SGD, momentum 0.9 & $10^{-2}$ & 0.7241 & 19 \\
SGD, momentum 0.5 & $10^{-2}$ & 0.7132 & 18 \\
SGD, momentum 0.99 & $10^{-2}$ & 0.6606 & 18 \\
\bottomrule
\end{tabular}
\end{table}

\reportfigure{0.94}{../results/q3/q3_curves.png}{Training and validation loss and accuracy for the main optimizer comparison.}
\reportfigure{0.94}{../results/q3/q3_momentum.png}{Validation behaviour across the tested SGD momentum values.}

\subsection{Custom architecture comparison}
Model A produced the higher test accuracy (79.82\% versus 76.18\%) and higher
macro precision and recall. Model B reduced trainable parameters by about
81\% and MACs by about 88\% relative to Model A. Its serialized model file
was about one quarter the size. Median epoch times were similar in the
recorded environment, so the principal observed advantage was lower model
and compute footprint rather than a large wall-clock speedup.

\begin{table}[H]\centering
\caption{Custom CNN performance and computational footprint.}\label{tab:q4}
\small
\begin{tabular}{lrr}
\toprule
Metric & Model A & Model B \\
\midrule
Trainable parameters & 422,437 & 80,102 \\
Non-trainable parameters & 960 & 1,414 \\
Total parameters & 423,397 & 81,516 \\
MACs per image & 60.20 M & 7.34 M \\
Model size on disk & 1.69 MiB & 0.42 MiB \\
Median time per epoch & 10.2 s & 9.3 s \\
Best validation accuracy & 79.85\% (epoch 21) & 74.59\% (epoch 25) \\
Test accuracy & 79.82\% & 76.18\% \\
Macro precision & 81.05\% & 75.99\% \\
Macro recall & 78.97\% & 75.73\% \\
\bottomrule
\end{tabular}
\end{table}

The saved class-level metrics show differences by category. Model A has its
highest recall on sunflowers (89.5\%) and dandelions (88.8\%), while daisy
recall is lowest at 65.3\%. Model B reaches 91.4\% recall on sunflowers but
64.6\% on roses. These per-class values suggest that aggregate accuracy
conceals different class-specific error patterns.

\reportfigure{0.82}{../results/q4/curves_A.png}{Model A training and validation curves.}
\reportfigure{0.82}{../results/q4/confusion_A.png}{Model A test confusion matrix.}
\reportfigure{0.82}{../results/q4/curves_B.png}{Model B training and validation curves.}
\reportfigure{0.82}{../results/q4/confusion_B.png}{Model B test confusion matrix.}

\subsection{Pretrained model comparison}
MobileNetV2 achieved the highest recorded test accuracy (85.09\%) and macro
precision (85.72\%) and recall (85.12\%). EfficientNet-B0 scored 83.82\% test
accuracy. MobileNetV2 also had fewer parameters, lower MACs, and a smaller
saved file than EfficientNet-B0 in this setup. Both pretrained models exceed
the custom Model B accuracy, but Model B remains substantially smaller and
less computationally expensive.

\begin{table}[H]\centering
\caption{Transfer-learning results on the held-out test split.}\label{tab:q5}
\small
\begin{tabular}{lrr}
\toprule
Metric & MobileNetV2 & EfficientNet-B0 \\
\midrule
Total parameters & 2,264,389 & 4,055,976 \\
Model size on disk & 9.20 MB & 16.28 MB \\
MACs per image & 97.8 M & 126.0 M \\
Input to backbone & $128\times128$ & $128\times128$ \\
Best validation accuracy & 86.03\% (epoch 16) & 84.57\% (epoch 20) \\
Test accuracy & 85.09\% & 83.82\% \\
Macro precision & 85.72\% & 84.07\% \\
Macro recall & 85.12\% & 83.44\% \\
Median epoch time, head stage & 18.7 s & 34.3 s \\
Median epoch time, fine-tune stage & 22.0 s & 36.6 s \\
\bottomrule
\end{tabular}
\end{table}

\reportfigure{0.82}{../results/q5/curves_mobilenetv2.png}{MobileNetV2 training and validation curves; the figure marks the fine-tuning stage and best validation epoch.}
\reportfigure{0.82}{../results/q5/confusion_mobilenetv2.png}{MobileNetV2 test confusion matrix.}
\reportfigure{0.82}{../results/q5/curves_efficientnetb0.png}{EfficientNet-B0 training and validation curves.}
\reportfigure{0.82}{../results/q5/confusion_efficientnetb0.png}{EfficientNet-B0 test confusion matrix.}

\section{Discussion}
The results support two practical observations. First, the standard custom CNN
was more accurate than the depthwise separable custom CNN, while the latter
used far fewer parameters and MACs. This is a useful trade-off when deployment
resources matter. Second, transfer learning improved test accuracy on this
dataset: MobileNetV2 gained 5.27 percentage points over Model A and 8.91
points over Model B. These gains came with approximately 27.8 times the
parameters and 13.3 times the MACs compared with Model B.

MobileNetV2 is the strongest option when maximizing the recorded accuracy is
the priority. Model B is more suitable when model size and inference cost are
the dominant constraints. EfficientNet-B0 did not surpass MobileNetV2 here and
required more parameters, MACs, model storage, and epoch time in the saved
experiment.

The comparison has limitations. Only one fixed stratified split and one set
of training runs are reported; no repeated-seed confidence intervals are
available. The timing values depend on the unrecorded execution hardware and
should be interpreted only as within-run comparisons. Upsampling a 64-pixel
image to 128 pixels does not add image detail. The repository also does not
record a per-class evaluation for the pretrained models in the summary tables,
although full confusion matrices and class metrics are saved. Finally, the
image resizing step may distort aspect ratios.

\section{Conclusion}
The project implemented a reproducible flower-classification pipeline and
compared optimizer choices, custom CNN architectures, and pretrained
backbones. Adam provided the best validation result in the optimizer sweep.
Among custom models, Model A achieved 79.82\% test accuracy, while Model B
reduced computation and parameter count substantially with a modest accuracy
drop. MobileNetV2 achieved the best recorded test accuracy at 85.09\%, at a
larger computational and storage cost. The results demonstrate that the
preferred model depends on whether predictive performance or lightweight
deployment is the main requirement.

\section*{Reproducibility and repository}
The source code, saved metric files, plots, and serialized models are included
in the project repository. The scripts are \class{prepare\_data.py},
\class{train\_optimizers.py}, \class{models.py}, \class{train\_eval.py}, and
\class{train\_sota.py}. The repository README documents the required Python
packages and run order. The results in this report are transcribed from
\class{results/q3}, \class{results/q4}, and \class{results/q5}.

\begin{thebibliography}{9}
\raggedright
\bibitem{tfdsflowers}
TensorFlow Datasets, ``tf\_flowers,'' dataset catalogue.
\url{https://www.tensorflow.org/datasets/catalog/tf_flowers}.

\bibitem{mobilenet}
M. Sandler, A. Howard, M. Zhu, A. Zhmoginov, and L.-C. Chen,
``MobileNetV2: Inverted Residuals and Linear Bottlenecks,''
\emph{Proceedings of the IEEE Conference on Computer Vision and Pattern
Recognition}, 2018.

\bibitem{efficientnet}
M. Tan and Q. V. Le, ``EfficientNet: Rethinking Model Scaling for
Convolutional Neural Networks,'' \emph{Proceedings of the 36th International
Conference on Machine Learning}, 2019.

\bibitem{uomemblem}
University of Moratuwa, ``The Emblem of University.''
\url{https://uom.lk/emblem-university.php}.
\end{thebibliography}

\end{document}
